\begin{tikzpicture}[
    font=\footnotesize,
    >=Stealth,
    % -------------------------------------------------------------------------
    % Node Styling
    % -------------------------------------------------------------------------
    base/.style={
        rectangle, 
        rounded corners=3pt, 
        thick, 
        align=center, 
        inner sep=3.5pt
    },
    root/.style={
        base,
        draw=blue!75!black,
        fill=blue!10,
        text width=9.8cm,
        minimum height=0.7cm,
        font=\small\bfseries
    },
    fixedbox/.style={
        base,
        draw=blue!75!black,
        fill=blue!8,
        text width=3.3cm,
        minimum height=0.65cm,
        font=\scriptsize\bfseries
    },
    fixeditem/.style={
        base,
        rounded corners=2pt,
        draw=blue!50!black,
        fill=blue!2,
        align=left,
        text width=3.3cm,
        inner sep=3.5pt,
        font=\tiny
    },
    variablebox/.style={
        base,
        draw=orange!75!black,
        fill=orange!8,
        text width=3.3cm,
        minimum height=0.65cm,
        font=\scriptsize\bfseries
    },
    modeparent/.style={
        base,
        rounded corners=3pt,
        text width=2.9cm,
        minimum height=0.65cm,
        font=\scriptsize\bfseries
    },
    subenv/.style={
        base,
        rounded corners=2pt,
        text width=1.35cm,
        minimum height=0.55cm,
        font=\scriptsize
    },
    stratbanner/.style={
        base,
        draw=purple!65!black,
        fill=purple!8,
        text width=7.4cm,
        minimum height=0.55cm,
        font=\scriptsize\bfseries
    },
    stratbox/.style={
        base,
        rounded corners=2pt,
        draw=purple!60!black,
        fill=purple!5,
        text width=1.7cm,
        minimum height=0.55cm,
        inner sep=2pt,
        font=\scriptsize
    },
    treearrow/.style={
        ->,
        >=Stealth,
        thick,
        draw=black!70
    }
]

    % =========================================================================
    % 1. ROOT
    % =========================================================================
    \node[root] (root) {
        Prompt Specification\\[1pt]
    };

    % =========================================================================
    % 2. TOP LEVEL (INVARIANT vs. FACTORIAL)
    % =========================================================================
    \node[
        fixedbox, 
        below=1.1cm of root.south,
        xshift=-3.9cm
    ] (fixed) {
        Invariant Prompts\\[1pt]
    };

    \node[
        variablebox, 
        below=1.1cm of root.south,
        xshift=1.9cm
    ] (variable) {
        Variant Prompts
    };

    \draw[treearrow] (root.south) -- ++(0,-0.4) -| (fixed.north);
    \draw[treearrow] (root.south) -- ++(0,-0.4) -| (variable.north);

    % =========================================================================
    % 3. LEFT COLUMN: BUS-STYLE ROUTED STACK
    % =========================================================================
    \node[fixeditem, below=0.85cm of fixed] (layout) {
        \textbf{Task Layout}\\
        Target identifier, dimension $D$, bounds $[-5, 5]^D$, evaluation budget $B_{\text{synth}}$, and tags.
    };

    \node[fixeditem, below=0.85cm of layout] (example) {
        \textbf{\texttt{EXAMPLE}}\\
        Common algorithm class, \texttt{\_\_init\_\_}, \texttt{\_\_call\_\_}, bound extraction, budget tracker.
    };

    \node[fixeditem, below=0.85cm of example] (format) {
        \textbf{\texttt{FORMAT}}\\
        Uniform response schema, single block, return tuple $(\mathbf{x}^*, y_{\text{raw}})$, whitelist NumPy.
    };

    % Clean, dedicated left bus line so arrows don't shoot straight through nodes
    \draw[treearrow, rounded corners=2pt] (fixed.west) -- ++(-0.35,0) |- (layout.west);
    \draw[treearrow, rounded corners=2pt] (fixed.west) -- ++(-0.35,0) |- (example.west);
    \draw[treearrow, rounded corners=2pt] (fixed.west) -- ++(-0.35,0) |- (format.west);

    % =========================================================================
    % 4. RIGHT COLUMN: IMPLICIT & EXPLICIT (CO-EQUAL LEVEL)
    % =========================================================================
    \node[
        modeparent,
        draw=yellow!65!black,
        fill=yellow!12,
        below=0.8cm of variable.south,
        xshift=-1.9cm
    ] (implicit) {
        \textbf{Implicit Mode}\\[1pt]
        \tiny Incomplete context; evaluation variability undisclosed
    };

    \node[
        modeparent,
        draw=orange!75!black,
        fill=orange!12,
        below=0.8cm of variable.south,
        xshift=1.9cm
    ] (explicit) {
        \textbf{Explicit Mode}\\[1pt]
        \tiny Evaluation behavior is explicitly characterized
    };

    \draw[treearrow] (variable.south) -- ++(0,-0.3) -| (implicit.north);
    \draw[treearrow] (variable.south) -- ++(0,-0.3) -| (explicit.north);

    % =========================================================================
    % 5. EXPLICIT SUB-BRANCHES: CLEAN & NOISY
    % =========================================================================
    \node[
        subenv,
        draw=green!60!black,
        fill=green!8,
        below=0.75cm of explicit.south,
        xshift=-0.85cm
    ] (clean) {
        \textbf{Clean}
    };

    \node[
        subenv,
        draw=red!60!black,
        fill=red!8,
        below=0.75cm of explicit.south,
        xshift=0.85cm
    ] (noisy) {
        \textbf{Noisy}
    };

    \draw[treearrow] (explicit.south) -- ++(0,-0.25) -| (clean.north);
    \draw[treearrow] (explicit.south) -- ++(0,-0.25) -| (noisy.north);

    % =========================================================================
    % 6. STRATEGY SCAFFOLD BANNER (CLEAR VERTICAL GAP BELOW CLEAN/NOISY)
    % =========================================================================
    \node[
        stratbanner,
        below=0.9cm of clean.south -| variable.south
    ] (strat_title) {
        \textbf{Strategy Scaffold}
    };

    % Direct vertical drops into the banner
    \draw[treearrow] (implicit.south) -- (implicit.south |- strat_title.north);
    \draw[treearrow] (clean.south)    -- (clean.south |- strat_title.north);
    \draw[treearrow] (noisy.south)    -- (noisy.south |- strat_title.north);

    % =========================================================================
    % 7. FOUR STRATEGY BLOCKS (GENEROUS BOTTOM MARGIN)
    % =========================================================================
    \node[stratbox, below=0.6cm of strat_title.south west, anchor=north west, xshift=0.08cm] (s_base) {
        \textbf{Baseline}
    };

    \node[stratbox, right=0.18cm of s_base] (s_think) {
        \textbf{Thinking}
    };

    \node[stratbox, right=0.18cm of s_think] (s_vec) {
        \textbf{Vectorization}
    };

    \node[stratbox, right=0.18cm of s_vec] (s_guided) {
        \textbf{Guided}
    };

    \draw[treearrow] (strat_title.south -| s_base.north)   -- (s_base.north);
    \draw[treearrow] (strat_title.south -| s_think.north)  -- (s_think.north);
    \draw[treearrow] (strat_title.south -| s_vec.north)    -- (s_vec.north);
    \draw[treearrow] (strat_title.south -| s_guided.north) -- (s_guided.north);

\end{tikzpicture}